\section{Выводы}
Выполнив первую лабораторную работу по курсу \enquote{Информационный поиск}, я
научился собирать корпус документов из публичных API. Основное, что я узнал
нового:

\begin{itemize}
    \item работа с реальными API (Last.fm, Lyrics.ovh, MusicBrainz): ключи API,
    параметры запросов, разбор ответов в формате JSON;
    \item обработка сбоев сети при загрузке большого объёма данных: таймауты,
    повторные попытки, ограничение частоты запросов --- без этого загрузка
    корпуса обрывается на первых же ошибках (что и произошло в первой версии
    программы);
    \item разбиение данных на документы и хранение метаданных отдельно от
    содержимого документа;
    \item подсчёт статистики корпуса: размер сырых данных, количество
    документов, размер выделенного текста, средние размеры.
\end{itemize}

Полученный корпус будет использован в следующих лабораторных работах для
построения поискового индекса и реализации поиска. Работа с существующими
поисковиками (Google с ограничением по сайтам, встроенный поиск MusicBrainz)
показала их недостатки --- дубликаты, нерелевантную выдачу и поиск только по
метаданным; именно эти недостатки предстоит устранить собственным поисковиком.

\pagebreak
